% --- Archon physics formalization source begin ---
% archon:physics
% archon:chemistry
% archon:covers IChO2026Problems/problem_icho_2026_t4_a5.lean
% archon:source-report reports/icho_2026/problem_icho_2026_t4_a5.source.json
% archon:problem-id icho_2026_t4
% archon:part-id T4-A5

\section{Icho Chemistry problem icho\_2026\_t4\_a5}

\paragraph{Problem source.}
\#\# Source
58th International Chemistry Olympiad, Tashkent, Uzbekistan, 2026.
Official-materials index: https://www.icho2026.uz/problems
English problem PDF: ../icho\_2026\_source/raw/theory\_problem.pdf
English solution/rubric PDF: ../icho\_2026\_source/raw/theory\_solution.pdf
Problem PDF page 38; printed local question page 2; solution starts on PDF page 35.
The page PNG is primary visual evidence for formulas, structures, tables, and figures.

\#\# T4. The Nuclear Past of Uzbekistan
T4. The nuclear past of Uzbekistan 6\% Natural uranium consists primarily of two isotopes, 235U and 238U . 235U is responsible for sustaining nuclear fission reactions. 4.1 Calculate the atomic abundance of 235U in natural uranium, assuming it consists only of isotopes 235U (235.04 a.u.) and 238U (238.05 a.u.). The main reaction occurring in nuclear power plants is the reaction where 235U absorbs a neutron and undergoes fission, releasing energy and producing three additional neutrons, thus sustaining a chain reaction. 235 92 U +1 0 n →… + 31 0n Neutrons that induce the chain reaction are produced by several ways: a) Reaction of 11B with 𝛼-particles inside the reactor core b) Interaction of 𝛾rays with 2D nuclei present in the reactor coolant water c) Reaction of 9Be with 𝛼-particles produced in an externally installed neutron source. 4.2 Write the nuclear reaction equations for (a), (b), and (c). The graph above displays the percentage yield, P(A), of fission products by mass number, A, resulting from 235U fission. 4.3 Write the most common nuclear fission equation for splitting of 235U upon neutron absorption. The pair of elements formed in this reaction are in the same group of Periodic Table. 4.4 Calculate the energy released in this reaction (ΔE, MeV), if the binding energy 𝐵𝐸(235U) = 7.59 MeV/nucleon and average binding energy for fission products 𝐵𝐸(fis.) = 8.45 MeV/nucleon. Neglect the binding energy of free neutrons. Fission neutrons are produced with an average energy of 2 MeV and slow down through many scatterings with nuclei. After several collisions, their speed decreases until their average kinetic energy matches that of the surrounding atoms or molecules. The logarithmic energy decrement, 𝜉, shows the average decrease per collision in the logarithm of the neutron energy: 𝜉= ln ( 𝐸𝑖𝑛𝑖𝑡𝑖𝑎𝑙 𝐸𝑓𝑖𝑛𝑎𝑙) 𝐸𝑖𝑛𝑖𝑡𝑖𝑎𝑙and 𝐸𝑓𝑖𝑛𝑎𝑙are the initial and final neutron energies per collision. 𝜉is a constant for each type of material and does not depend on the initial neutron energy.

\#\# Current subquestion 4.5
Determine the average number of collisions required, n𝑐, to slow a neutron from 2 MeV to 0.012 eV, using water as the moderator. 𝜉(water)=0.948

\#\# Dataset metadata
IChO 2026; T4; raw rubric points=2; kind=theory; formalization\_ready=true.

\paragraph{Shared problem context.}
T4. The nuclear past of Uzbekistan 6\% Natural uranium consists primarily of two isotopes, 235U and 238U . 235U is responsible for sustaining nuclear fission reactions. 4.1 Calculate the atomic abundance of 235U in natural uranium, assuming it consists only of isotopes 235U (235.04 a.u.) and 238U (238.05 a.u.). The main reaction occurring in nuclear power plants is the reaction where 235U absorbs a neutron and undergoes fission, releasing energy and producing three additional neutrons, thus sustaining a chain reaction. 235 92 U +1 0 n →… + 31 0n Neutrons that induce the chain reaction are produced by several ways: a) Reaction of 11B with 𝛼-particles inside the reactor core b) Interaction of 𝛾rays with 2D nuclei present in the reactor coolant water c) Reaction of 9Be with 𝛼-particles produced in an externally installed neutron source. 4.2 Write the nuclear reaction equations for (a), (b), and (c). The graph above displays the percentage yield, P(A), of fission products by mass number, A, resulting from 235U fission. 4.3 Write the most common nuclear fission equation for splitting of 235U upon neutron absorption. The pair of elements formed in this reaction are in the same group of Periodic Table. 4.4 Calculate the energy released in this reaction (ΔE, MeV), if the binding energy 𝐵𝐸(235U) = 7.59 MeV/nucleon and average binding energy for fission products 𝐵𝐸(fis.) = 8.45 MeV/nucleon. Neglect the binding energy of free neutrons. Fission neutrons are produced with an average energy of 2 MeV and slow down through many scatterings with nuclei. After several collisions, their speed decreases until their average kinetic energy matches that of the surrounding atoms or molecules. The logarithmic energy decrement, 𝜉, shows the average decrease per collision in the logarithm of the neutron energy: 𝜉= ln ( 𝐸𝑖𝑛𝑖𝑡𝑖𝑎𝑙 𝐸𝑓𝑖𝑛𝑎𝑙) 𝐸𝑖𝑛𝑖𝑡𝑖𝑎𝑙and 𝐸𝑓𝑖𝑛𝑎𝑙are the initial and final neutron energies per collision. 𝜉is a constant for each type of material and does not depend on the initial neutron energy.

\paragraph{Current subquestion.}
Determine the average number of collisions required, n𝑐, to slow a neutron from 2 MeV to 0.012 eV, using water as the moderator. 𝜉(water)=0.948

\paragraph{Recorded answer/context.}
𝑛𝑐= ln( 𝐸𝑖𝑛𝑖𝑡𝑖𝑎𝑙 𝐸𝑓𝑖𝑛𝑎𝑙) 𝜉 = 19.97 ≈20 2 points for the correct answer 20. 1 point may be given for the right formula but the wrong answer. 2.0 pt

\paragraph{Figure/image paths.}
\begin{itemize}
\item ../icho\_2026\_source/image/T4\_page-2.png
\item ../icho\_2026\_source/image/T4\_page-1.png
\end{itemize}

\subsection{Formalization contract and informal proof}

The requested collision count is an average and is therefore a nonnegative
real number, not an integer-valued count.  Energies are represented by their
numerical values in electronvolts; thus the source's $2\,\mathrm{MeV}$ is
$2{,}000{,}000\,\mathrm{eV}$.  The reported answer is a nearest-integer
presentation of that average.

\begin{definition}

\leanok
[Electronvolt energy readout]
\label{def:physics:icho_2026_t4_a5:energy_ev}
\lean{IChO2026Problems.T4A5.EnergyEV}
An energy is represented by its real numerical value in electronvolts.
\end{definition}
\begin{proof}

\leanok
The source data are converted to and manipulated in the common eV scale.
\end{proof}

\begin{definition}

\leanok
[Average collision-count readout]
\label{def:physics:icho_2026_t4_a5:average_collision_count}
\lean{IChO2026Problems.T4A5.AverageCollisionCount}
An average collision count is represented by a nonnegative real readout.
\end{definition}
\begin{proof}

\leanok
The physical quantity is an average, so its formal representation need not be
an integer.
\end{proof}

\begin{definition}

\leanok
[Moderator identity]
\label{def:physics:icho_2026_t4_a5:moderator}
\lean{IChO2026Problems.T4A5.Moderator}
The local moderator identity distinguishes water, the material supplied in
T4-A5.
\end{definition}
\begin{proof}

\leanok
The identity records which material's logarithmic decrement is being used.
\end{proof}

\begin{definition}\leanok
[Neutron moderation data]
\label{def:physics:icho_2026_t4_a5:data}
\lean{IChO2026Problems.T4A5.NeutronModerationData}
\uses{def:physics:icho_2026_t4_a5:energy_ev, def:physics:icho_2026_t4_a5:average_collision_count, def:physics:icho_2026_t4_a5:moderator}
The data record contains the positive initial and final neutron energies in
eV, the positive moderator logarithmic energy decrement $\xi$, and the
nonnegative average collision count $n_c$.
\end{definition}
\begin{proof}

\leanok
These are the physical quantities named in T4-A5.  Positivity ensures that
the energy ratio and its logarithm are meaningful and that division by
$\xi$ is valid.
\end{proof}

\begin{definition}\leanok
[Cumulative logarithmic-decrement law]
\label{def:physics:icho_2026_t4_a5:moderation_law}
\lean{IChO2026Problems.T4A5.NeutronModerationData.satisfiesModerationLaw}
\uses{def:physics:icho_2026_t4_a5:data}
The moderator law states
\[
 n_c\,\xi=\log(E_{\mathrm{initial}}/E_{\mathrm{final}}).
\]
It accumulates the source's average logarithmic decrease per collision over
the average number of collisions; the decrement is independent of initial
energy as stipulated in the question.
\end{definition}
\begin{proof}

\leanok
Each collision contributes the same mean logarithmic decrement.  Summing
those $n_c$ contributions gives the logarithm of the total energy ratio.
\end{proof}

\begin{definition}\leanok
[Water moderation inputs]
\label{def:physics:icho_2026_t4_a5:water_data}
\lean{IChO2026Problems.T4A5.NeutronModerationData.matchesWaterData}
\uses{def:physics:icho_2026_t4_a5:data}
For water, the supplied values are
\[
 E_{\mathrm{initial}}=2{,}000{,}000\ \mathrm{eV},\qquad
 E_{\mathrm{final}}=0.012\ \mathrm{eV},\qquad \xi=0.948.
\]
\end{definition}
\begin{proof}

\leanok
The first equality is the stated MeV-to-eV conversion; the latter two are the
given final energy and water decrement.  None is a collision-count answer.
\end{proof}

\begin{lemma}\leanok
[Collision-count formula]
\label{lem:physics:icho_2026_t4_a5:collision_count_formula}
\lean{IChO2026Problems.T4A5.collision_count_formula}
\uses{def:physics:icho_2026_t4_a5:data, def:physics:icho_2026_t4_a5:moderation_law}
Every positive moderation datum satisfying the cumulative law has
\[
 n_c=\frac{\log(E_{\mathrm{initial}}/E_{\mathrm{final}})}{\xi}.
\]
\end{lemma}
\begin{proof}

\leanok
The law is a linear equation in $n_c$.  Since $\xi>0$, division by $\xi$
uniquely isolates the average collision count.
\end{proof}

\begin{theorem}\leanok
[Water collision count]
\label{thm:physics:icho_2026_t4_a5:water_collision_count}
\lean{IChO2026Problems.T4A5.water_collision_count}
\uses{def:physics:icho_2026_t4_a5:water_data, lem:physics:icho_2026_t4_a5:collision_count_formula}
For water, the exact model readout is
\[
 n_c=\frac{\log(2{,}000{,}000/0.012)}{0.948},
\]
and the source's $19.97\approx20$ report is encoded by
$|n_c-19.97|<10^{-4}$ together with $|n_c-20|<\tfrac12$.
\end{theorem}
\begin{proof}

\leanok
Substitute the three water data equations into the collision-count formula.
The resulting real number is approximately $19.9699434$.  Explicit logarithm
bounds put it within $10^{-4}$ of the reported two-decimal value $19.97$ and
within one half of the reported nearest integer $20$; neither rounded value is
treated as an exact equality.
\end{proof}

\begin{theorem}\leanok
[IChO T4-A5 target]
\label{thm:physics:icho_2026_t4_a5:target}
\lean{IChO2026Problems.T4A5.average_collisions_for_water}
\uses{thm:physics:icho_2026_t4_a5:water_collision_count}
The covered Lean file presents the source formula, its two-decimal value
$19.97$, and its nearest-integer water result $20$, with the total
logarithmic-decrement law as the load-bearing governing assumption.
\end{theorem}
\begin{proof}

\leanok
This is the requested conclusion of the preceding water calculation.
\end{proof}
% --- Archon physics formalization source end ---
